\begin{beginnerstatement}[Kurz erklärt: ein Projekt erzeugen]

Das Kommando \texttt{init} startet den Generator für ein neues
Produktprojekt. Der Generator baut aus dem gewählten Profil ein
\emph{Scaffold}, also ein vorbereitetes Grundgerüst aus Ordnern und Dateien.
Das Zielverzeichnis ist der Ort, an dem dieses Projekt angelegt werden soll.
Ein \emph{Dry Run} prüft den vorgesehenen Ablauf, ohne Dateien zu schreiben. Zur
Projektidentität gehören Angaben wie Name und Kurzname, durch die das erzeugte
Projekt vom Template unterschieden wird. Eine Anwendungs-ID kennzeichnet eine
Tauri-Anwendung eindeutig. Ein \emph{Reverse-Domain-Identifier} ist eine solche
ID, bei der Wörter wie in \texttt{com.example.product} von der allgemeinen
Internet-Domain zum konkreten Produkt angeordnet sind.

\end{beginnerstatement}
